\section{成化正德的冲突、灾害与司法证据}

战争、灾荒、疾病、赈恤、案狱和边地暴力并不是彼此独立的材料领域。军报和判牍往往由胜者、官府或审讯机关形成；灾情、医疗与慈善记录也常以请赈、捐输、工程或死亡数字的形式出现。它们能显示国家和地方组织如何命名风险与行动，却不能把受害者、病者、被征发者或被定罪者的经验全部还原。

\subsection{军政冲突与地方动员}

\paragraph{1481（成化十七年）：田土、赋役、盐课或漕运的本年政策记录提供财政执行的候选节点，需同地区册籍比照。（材料核查重点：报称信息的来源、抄传与行政分类）} 此一账本锚点将冲突、风险或法律处置置于具体时段。它所能可靠支持的判断是来源导向的年度桥接条目：本年材料可定位相关政务线索；具体月日、发文机关、地域和执行结果仍须逐项回查，不能以制度文本或奏报替代实际效果。材料链由，因而同时保存了命令、报送、传记、地方记忆或跨政权叙事的不同立场。问题形成的条件是财政征解、交通工程或货币支付的压力使相关措施进入政务；随后产生的可见影响是其法定安排不等于地方执行，后续须以地域材料追踪负担与成效。对于医疗救助、慈善网络、普通人的伤亡与案狱处境，材料往往尤为稀薄。桥接条目只标示可回查的年度问题与材料链；实录有中央选择性，崇祯期尤其需要奏疏、地方、朝鲜及满文材料交叉。；本条特别核对报称信息的来源、抄传与行政分类。

\paragraph{1481（成化十七年）：本年灾荒、河工或赈恤的报奏材料可与地方志互校，报灾范围和实际损失应分别记录。（材料核查重点：诏令或奏议的形成与发受链）} 此一账本锚点将冲突、风险或法律处置置于具体时段。它所能可靠支持的判断是来源导向的年度桥接条目：本年材料可定位相关政务线索；具体月日、发文机关、地域和执行结果仍须逐项回查，不能以制度文本或奏报替代实际效果。材料链由，因而同时保存了命令、报送、传记、地方记忆或跨政权叙事的不同立场。问题形成的条件是自然风险与地方报告促成朝廷或地方的应对记录；随后产生的可见影响是赈济、迁徙和损失的实际范围仍须以受灾地区材料分别核验。对于医疗救助、慈善网络、普通人的伤亡与案狱处境，材料往往尤为稀薄。桥接条目只标示可回查的年度问题与材料链；实录有中央选择性，崇祯期尤其需要奏疏、地方、朝鲜及满文材料交叉。；本条特别核对诏令或奏议的形成与发受链。

\paragraph{1481（成化十七年）：本年灾荒、河工或赈恤的报奏材料可与地方志互校，报灾范围和实际损失应分别记录。（材料核查重点：报称信息的来源、抄传与行政分类）} 此一账本锚点将冲突、风险或法律处置置于具体时段。它所能可靠支持的判断是来源导向的年度桥接条目：本年材料可定位相关政务线索；具体月日、发文机关、地域和执行结果仍须逐项回查，不能以制度文本或奏报替代实际效果。材料链由，因而同时保存了命令、报送、传记、地方记忆或跨政权叙事的不同立场。问题形成的条件是自然风险与地方报告促成朝廷或地方的应对记录；随后产生的可见影响是赈济、迁徙和损失的实际范围仍须以受灾地区材料分别核验。对于医疗救助、慈善网络、普通人的伤亡与案狱处境，材料往往尤为稀薄。桥接条目只标示可回查的年度问题与材料链；实录有中央选择性，崇祯期尤其需要奏疏、地方、朝鲜及满文材料交叉。；本条特别核对报称信息的来源、抄传与行政分类。

\paragraph{1482（成化十八年）：哈密、吐鲁番及周边朝贡关系的本年交涉显示东天山政治流动，册封不等于稳定控制} 此一账本锚点将冲突、风险或法律处置置于具体时段。它所能可靠支持的判断是来源导向的年度桥接条目：本年材料可定位相关政务线索；具体月日、发文机关、地域和执行结果仍须逐项回查，不能以制度文本或奏报替代实际效果。材料链由，因而同时保存了命令、报送、传记、地方记忆或跨政权叙事的不同立场。问题形成的条件是朝贡名义、边境安全与港口贸易的利益使交涉持续发生；随后产生的可见影响是它为后续册封、互市或海防安排留下文书与跨政权比较线索。对于医疗救助、慈善网络、普通人的伤亡与案狱处境，材料往往尤为稀薄。桥接条目只标示可回查的年度问题与材料链；实录有中央选择性，崇祯期尤其需要奏疏、地方、朝鲜及满文材料交叉。

\paragraph{1482（成化十八年）：成化、弘治间中枢任免、厂卫与言路的本年材料标记皇权和官僚关系的调整。（材料核查重点：诏令或奏议的形成与发受链）} 此一账本锚点将冲突、风险或法律处置置于具体时段。它所能可靠支持的判断是来源导向的年度桥接条目：本年材料可定位相关政务线索；具体月日、发文机关、地域和执行结果仍须逐项回查，不能以制度文本或奏报替代实际效果。材料链由，因而同时保存了命令、报送、传记、地方记忆或跨政权叙事的不同立场。问题形成的条件是继承、官僚协调或皇权运作的压力使问题进入朝廷程序；随后产生的可见影响是它影响后续人事与政策议程；动机和派系标签不能由单一叙事裁定。对于医疗救助、慈善网络、普通人的伤亡与案狱处境，材料往往尤为稀薄。桥接条目只标示可回查的年度问题与材料链；实录有中央选择性，崇祯期尤其需要奏疏、地方、朝鲜及满文材料交叉。；本条特别核对诏令或奏议的形成与发受链。

\paragraph{1482（成化十八年）：成化、弘治间中枢任免、厂卫与言路的本年材料标记皇权和官僚关系的调整。（材料核查重点：报称信息的来源、抄传与行政分类）} 此一账本锚点将冲突、风险或法律处置置于具体时段。它所能可靠支持的判断是来源导向的年度桥接条目：本年材料可定位相关政务线索；具体月日、发文机关、地域和执行结果仍须逐项回查，不能以制度文本或奏报替代实际效果。材料链由，因而同时保存了命令、报送、传记、地方记忆或跨政权叙事的不同立场。问题形成的条件是继承、官僚协调或皇权运作的压力使问题进入朝廷程序；随后产生的可见影响是它影响后续人事与政策议程；动机和派系标签不能由单一叙事裁定。对于医疗救助、慈善网络、普通人的伤亡与案狱处境，材料往往尤为稀薄。桥接条目只标示可回查的年度问题与材料链；实录有中央选择性，崇祯期尤其需要奏疏、地方、朝鲜及满文材料交叉。；本条特别核对报称信息的来源、抄传与行政分类。

\subsection{灾害、医疗与赈恤}

\paragraph{1482（成化十八年）：哈密王去世后，明廷围绕其继嗣与册封处理新的王位更替} 此一账本锚点将冲突、风险或法律处置置于具体时段。它所能可靠支持的判断是王位更替和册封程序可由实录与后出西域传互校；人名译写、亲属关系和绿洲各派支持范围在汉文材料中有异文。材料链由，因而同时保存了命令、报送、传记、地方记忆或跨政权叙事的不同立场。问题形成的条件是哈密王权既依赖本地王族联盟，也依赖明廷赏赐和吐鲁番压力下的外交平衡，继承空缺使各方都有介入机会；随后产生的可见影响是新王的名义册封维持了朝贡形式，却未解决吐鲁番的军事与贸易竞争；哈密危机在弘治朝继续反复。对于医疗救助、慈善网络、普通人的伤亡与案狱处境，材料往往尤为稀薄。仍须把战果、伤亡、执行和受害者的处境逐项分开核验。

\paragraph{1483（成化十九年）：哈密、吐鲁番及周边朝贡关系的本年交涉显示东天山政治流动，册封不等于稳定控制} 此一账本锚点将冲突、风险或法律处置置于具体时段。它所能可靠支持的判断是来源导向的年度桥接条目：本年材料可定位相关政务线索；具体月日、发文机关、地域和执行结果仍须逐项回查，不能以制度文本或奏报替代实际效果。材料链由，因而同时保存了命令、报送、传记、地方记忆或跨政权叙事的不同立场。问题形成的条件是朝贡名义、边境安全与港口贸易的利益使交涉持续发生；随后产生的可见影响是它为后续册封、互市或海防安排留下文书与跨政权比较线索。对于医疗救助、慈善网络、普通人的伤亡与案狱处境，材料往往尤为稀薄。桥接条目只标示可回查的年度问题与材料链；实录有中央选择性，崇祯期尤其需要奏疏、地方、朝鲜及满文材料交叉。

\paragraph{1483（成化十九年）：学校、考试、刻书或礼制材料在本年留下文化制度线索，精英文本不能代表全体社会} 此一账本锚点将冲突、风险或法律处置置于具体时段。它所能可靠支持的判断是来源导向的年度桥接条目：本年材料可定位相关政务线索；具体月日、发文机关、地域和执行结果仍须逐项回查，不能以制度文本或奏报替代实际效果。材料链由，因而同时保存了命令、报送、传记、地方记忆或跨政权叙事的不同立场。问题形成的条件是制度、教育或知识传播的需要推动了相关行动或文本形成；随后产生的可见影响是它提供制度和传播史的锚点，不能据此推断社会整体接受。对于医疗救助、慈善网络、普通人的伤亡与案狱处境，材料往往尤为稀薄。桥接条目只标示可回查的年度问题与材料链；实录有中央选择性，崇祯期尤其需要奏疏、地方、朝鲜及满文材料交叉。

\paragraph{1483（成化十九年）：荆襄、广西、辽东或西北的兵事和安置文书构成本年地方秩序重组的线索。（材料核查重点：诏令或奏议的形成与发受链）} 此一账本锚点将冲突、风险或法律处置置于具体时段。它所能可靠支持的判断是来源导向的年度桥接条目：本年材料可定位相关政务线索；具体月日、发文机关、地域和执行结果仍须逐项回查，不能以制度文本或奏报替代实际效果。材料链由，因而同时保存了命令、报送、传记、地方记忆或跨政权叙事的不同立场。问题形成的条件是前期边防、地方武装或跨区域资源调动的压力推动了该行动；随后产生的可见影响是它改变了后续调兵、补给或地方秩序的条件；战果和伤亡须分开核对。对于医疗救助、慈善网络、普通人的伤亡与案狱处境，材料往往尤为稀薄。桥接条目只标示可回查的年度问题与材料链；实录有中央选择性，崇祯期尤其需要奏疏、地方、朝鲜及满文材料交叉。；本条特别核对诏令或奏议的形成与发受链。

\paragraph{1483（成化十九年）：荆襄、广西、辽东或西北的兵事和安置文书构成本年地方秩序重组的线索。（材料核查重点：报称信息的来源、抄传与行政分类）} 此一账本锚点将冲突、风险或法律处置置于具体时段。它所能可靠支持的判断是来源导向的年度桥接条目：本年材料可定位相关政务线索；具体月日、发文机关、地域和执行结果仍须逐项回查，不能以制度文本或奏报替代实际效果。材料链由，因而同时保存了命令、报送、传记、地方记忆或跨政权叙事的不同立场。问题形成的条件是前期边防、地方武装或跨区域资源调动的压力推动了该行动；随后产生的可见影响是它改变了后续调兵、补给或地方秩序的条件；战果和伤亡须分开核对。对于医疗救助、慈善网络、普通人的伤亡与案狱处境，材料往往尤为稀薄。桥接条目只标示可回查的年度问题与材料链；实录有中央选择性，崇祯期尤其需要奏疏、地方、朝鲜及满文材料交叉。；本条特别核对报称信息的来源、抄传与行政分类。

\paragraph{1484（成化二十年）：学校、考试、刻书或礼制材料在本年留下文化制度线索，精英文本不能代表全体社会} 此一账本锚点将冲突、风险或法律处置置于具体时段。它所能可靠支持的判断是来源导向的年度桥接条目：本年材料可定位相关政务线索；具体月日、发文机关、地域和执行结果仍须逐项回查，不能以制度文本或奏报替代实际效果。材料链由，因而同时保存了命令、报送、传记、地方记忆或跨政权叙事的不同立场。问题形成的条件是制度、教育或知识传播的需要推动了相关行动或文本形成；随后产生的可见影响是它提供制度和传播史的锚点，不能据此推断社会整体接受。对于医疗救助、慈善网络、普通人的伤亡与案狱处境，材料往往尤为稀薄。桥接条目只标示可回查的年度问题与材料链；实录有中央选择性，崇祯期尤其需要奏疏、地方、朝鲜及满文材料交叉。

\subsection{审案、处分与法律语言}

\paragraph{1484（成化二十年）：本年灾荒、河工或赈恤的报奏材料可与地方志互校，报灾范围和实际损失应分别记录} 此一账本锚点将冲突、风险或法律处置置于具体时段。它所能可靠支持的判断是来源导向的年度桥接条目：本年材料可定位相关政务线索；具体月日、发文机关、地域和执行结果仍须逐项回查，不能以制度文本或奏报替代实际效果。材料链由，因而同时保存了命令、报送、传记、地方记忆或跨政权叙事的不同立场。问题形成的条件是自然风险与地方报告促成朝廷或地方的应对记录；随后产生的可见影响是赈济、迁徙和损失的实际范围仍须以受灾地区材料分别核验。对于医疗救助、慈善网络、普通人的伤亡与案狱处境，材料往往尤为稀薄。桥接条目只标示可回查的年度问题与材料链；实录有中央选择性，崇祯期尤其需要奏疏、地方、朝鲜及满文材料交叉。

\paragraph{1484（成化二十年）：田土、赋役、盐课或漕运的本年政策记录提供财政执行的候选节点，需同地区册籍比照。（材料核查重点：诏令或奏议的形成与发受链）} 此一账本锚点将冲突、风险或法律处置置于具体时段。它所能可靠支持的判断是来源导向的年度桥接条目：本年材料可定位相关政务线索；具体月日、发文机关、地域和执行结果仍须逐项回查，不能以制度文本或奏报替代实际效果。材料链由，因而同时保存了命令、报送、传记、地方记忆或跨政权叙事的不同立场。问题形成的条件是财政征解、交通工程或货币支付的压力使相关措施进入政务；随后产生的可见影响是其法定安排不等于地方执行，后续须以地域材料追踪负担与成效。对于医疗救助、慈善网络、普通人的伤亡与案狱处境，材料往往尤为稀薄。桥接条目只标示可回查的年度问题与材料链；实录有中央选择性，崇祯期尤其需要奏疏、地方、朝鲜及满文材料交叉。；本条特别核对诏令或奏议的形成与发受链。

\paragraph{1484（成化二十年）：田土、赋役、盐课或漕运的本年政策记录提供财政执行的候选节点，需同地区册籍比照。（材料核查重点：报称信息的来源、抄传与行政分类）} 此一账本锚点将冲突、风险或法律处置置于具体时段。它所能可靠支持的判断是来源导向的年度桥接条目：本年材料可定位相关政务线索；具体月日、发文机关、地域和执行结果仍须逐项回查，不能以制度文本或奏报替代实际效果。材料链由，因而同时保存了命令、报送、传记、地方记忆或跨政权叙事的不同立场。问题形成的条件是财政征解、交通工程或货币支付的压力使相关措施进入政务；随后产生的可见影响是其法定安排不等于地方执行，后续须以地域材料追踪负担与成效。对于医疗救助、慈善网络、普通人的伤亡与案狱处境，材料往往尤为稀薄。桥接条目只标示可回查的年度问题与材料链；实录有中央选择性，崇祯期尤其需要奏疏、地方、朝鲜及满文材料交叉。；本条特别核对报称信息的来源、抄传与行政分类。

\paragraph{1485（成化二十一年）：太监数量过万} 此一账本锚点将冲突、风险或法律处置置于具体时段。它所能可靠支持的判断是编年候选的年序可由官修与后出材料交叉；具体月日、参与者、战果或数字必须回查所列材料，不能将单一叙事当作定论。材料链由，因而同时保存了命令、报送、传记、地方记忆或跨政权叙事的不同立场。问题形成的条件是继承、官僚协调或皇权运作的压力使问题进入朝廷程序；随后产生的可见影响是它影响后续人事与政策议程；动机和派系标签不能由单一叙事裁定。对于医疗救助、慈善网络、普通人的伤亡与案狱处境，材料往往尤为稀薄。译写候选以编年材料定位；实录记载反映朝廷选择，崇祯期须另以奏疏、地方及敌对政权材料交叉。

\paragraph{1485（成化二十一年）：成化、弘治间中枢任免、厂卫与言路的本年材料标记皇权和官僚关系的调整} 此一账本锚点将冲突、风险或法律处置置于具体时段。它所能可靠支持的判断是来源导向的年度桥接条目：本年材料可定位相关政务线索；具体月日、发文机关、地域和执行结果仍须逐项回查，不能以制度文本或奏报替代实际效果。材料链由，因而同时保存了命令、报送、传记、地方记忆或跨政权叙事的不同立场。问题形成的条件是继承、官僚协调或皇权运作的压力使问题进入朝廷程序；随后产生的可见影响是它影响后续人事与政策议程；动机和派系标签不能由单一叙事裁定。对于医疗救助、慈善网络、普通人的伤亡与案狱处境，材料往往尤为稀薄。桥接条目只标示可回查的年度问题与材料链；实录有中央选择性，崇祯期尤其需要奏疏、地方、朝鲜及满文材料交叉。

\paragraph{1485（成化二十一年）：哈密、吐鲁番及周边朝贡关系的本年交涉显示东天山政治流动，册封不等于稳定控制。（材料核查重点：诏令或奏议的形成与发受链）} 此一账本锚点将冲突、风险或法律处置置于具体时段。它所能可靠支持的判断是来源导向的年度桥接条目：本年材料可定位相关政务线索；具体月日、发文机关、地域和执行结果仍须逐项回查，不能以制度文本或奏报替代实际效果。材料链由，因而同时保存了命令、报送、传记、地方记忆或跨政权叙事的不同立场。问题形成的条件是朝贡名义、边境安全与港口贸易的利益使交涉持续发生；随后产生的可见影响是它为后续册封、互市或海防安排留下文书与跨政权比较线索。对于医疗救助、慈善网络、普通人的伤亡与案狱处境，材料往往尤为稀薄。桥接条目只标示可回查的年度问题与材料链；实录有中央选择性，崇祯期尤其需要奏疏、地方、朝鲜及满文材料交叉。；本条特别核对诏令或奏议的形成与发受链。

\subsection{边地暴力与行政后果}

\paragraph{1485（成化二十一年）：哈密、吐鲁番及周边朝贡关系的本年交涉显示东天山政治流动，册封不等于稳定控制。（材料核查重点：报称信息的来源、抄传与行政分类）} 此一账本锚点将冲突、风险或法律处置置于具体时段。它所能可靠支持的判断是来源导向的年度桥接条目：本年材料可定位相关政务线索；具体月日、发文机关、地域和执行结果仍须逐项回查，不能以制度文本或奏报替代实际效果。材料链由，因而同时保存了命令、报送、传记、地方记忆或跨政权叙事的不同立场。问题形成的条件是朝贡名义、边境安全与港口贸易的利益使交涉持续发生；随后产生的可见影响是它为后续册封、互市或海防安排留下文书与跨政权比较线索。对于医疗救助、慈善网络、普通人的伤亡与案狱处境，材料往往尤为稀薄。桥接条目只标示可回查的年度问题与材料链；实录有中央选择性，崇祯期尤其需要奏疏、地方、朝鲜及满文材料交叉。；本条特别核对报称信息的来源、抄传与行政分类。

\paragraph{1486（成化二十二年）：本年灾荒、河工或赈恤的报奏材料可与地方志互校，报灾范围和实际损失应分别记录} 此一账本锚点将冲突、风险或法律处置置于具体时段。它所能可靠支持的判断是来源导向的年度桥接条目：本年材料可定位相关政务线索；具体月日、发文机关、地域和执行结果仍须逐项回查，不能以制度文本或奏报替代实际效果。材料链由，因而同时保存了命令、报送、传记、地方记忆或跨政权叙事的不同立场。问题形成的条件是自然风险与地方报告促成朝廷或地方的应对记录；随后产生的可见影响是赈济、迁徙和损失的实际范围仍须以受灾地区材料分别核验。对于医疗救助、慈善网络、普通人的伤亡与案狱处境，材料往往尤为稀薄。桥接条目只标示可回查的年度问题与材料链；实录有中央选择性，崇祯期尤其需要奏疏、地方、朝鲜及满文材料交叉。

\paragraph{1486（成化二十二年）：荆襄、广西、辽东或西北的兵事和安置文书构成本年地方秩序重组的线索} 此一账本锚点将冲突、风险或法律处置置于具体时段。它所能可靠支持的判断是来源导向的年度桥接条目：本年材料可定位相关政务线索；具体月日、发文机关、地域和执行结果仍须逐项回查，不能以制度文本或奏报替代实际效果。材料链由，因而同时保存了命令、报送、传记、地方记忆或跨政权叙事的不同立场。问题形成的条件是前期边防、地方武装或跨区域资源调动的压力推动了该行动；随后产生的可见影响是它改变了后续调兵、补给或地方秩序的条件；战果和伤亡须分开核对。对于医疗救助、慈善网络、普通人的伤亡与案狱处境，材料往往尤为稀薄。桥接条目只标示可回查的年度问题与材料链；实录有中央选择性，崇祯期尤其需要奏疏、地方、朝鲜及满文材料交叉。

\paragraph{1486（成化二十二年）：学校、考试、刻书或礼制材料在本年留下文化制度线索，精英文本不能代表全体社会。（材料核查重点：诏令或奏议的形成与发受链）} 此一账本锚点将冲突、风险或法律处置置于具体时段。它所能可靠支持的判断是来源导向的年度桥接条目：本年材料可定位相关政务线索；具体月日、发文机关、地域和执行结果仍须逐项回查，不能以制度文本或奏报替代实际效果。材料链由，因而同时保存了命令、报送、传记、地方记忆或跨政权叙事的不同立场。问题形成的条件是制度、教育或知识传播的需要推动了相关行动或文本形成；随后产生的可见影响是它提供制度和传播史的锚点，不能据此推断社会整体接受。对于医疗救助、慈善网络、普通人的伤亡与案狱处境，材料往往尤为稀薄。桥接条目只标示可回查的年度问题与材料链；实录有中央选择性，崇祯期尤其需要奏疏、地方、朝鲜及满文材料交叉。；本条特别核对诏令或奏议的形成与发受链。

\paragraph{1486（成化二十二年）：学校、考试、刻书或礼制材料在本年留下文化制度线索，精英文本不能代表全体社会。（材料核查重点：报称信息的来源、抄传与行政分类）} 此一账本锚点将冲突、风险或法律处置置于具体时段。它所能可靠支持的判断是来源导向的年度桥接条目：本年材料可定位相关政务线索；具体月日、发文机关、地域和执行结果仍须逐项回查，不能以制度文本或奏报替代实际效果。材料链由，因而同时保存了命令、报送、传记、地方记忆或跨政权叙事的不同立场。问题形成的条件是制度、教育或知识传播的需要推动了相关行动或文本形成；随后产生的可见影响是它提供制度和传播史的锚点，不能据此推断社会整体接受。对于医疗救助、慈善网络、普通人的伤亡与案狱处境，材料往往尤为稀薄。桥接条目只标示可回查的年度问题与材料链；实录有中央选择性，崇祯期尤其需要奏疏、地方、朝鲜及满文材料交叉。；本条特别核对报称信息的来源、抄传与行政分类。

\paragraph{1487（成化二十三年）九月1日：成化皇帝病重} 此一账本锚点将冲突、风险或法律处置置于具体时段。它所能可靠支持的判断是编年候选的年序可由官修与后出材料交叉；具体月日、参与者、战果或数字必须回查所列材料，不能将单一叙事当作定论。材料链由，因而同时保存了命令、报送、传记、地方记忆或跨政权叙事的不同立场。问题形成的条件是继承、官僚协调或皇权运作的压力使问题进入朝廷程序；随后产生的可见影响是它影响后续人事与政策议程；动机和派系标签不能由单一叙事裁定。对于医疗救助、慈善网络、普通人的伤亡与案狱处境，材料往往尤为稀薄。译写候选以编年材料定位；实录记载反映朝廷选择，崇祯期须另以奏疏、地方及敌对政权材料交叉。
